\documentclass[10pt,a4paper]{amsart}
\usepackage[margin=19mm]{geometry}
\usepackage{amsmath,amssymb,mathtools}
\usepackage[scr=rsfs,cal=euler]{mathalfa}
\usepackage{xfrac}
\usepackage[unicode,hidelinks,bookmarks=false]{hyperref}
\newcommand{\ie}{i.e.\ }
\newcommand{\cf}{cf.\ }
\newcommand{\resp}{respectively\ }
\newcommand{\Subsection}{\subsection}
\providecommand{\nobreakdash}{\nobreak\hbox{-}}
\setlength{\emergencystretch}{2em}
\multlinegap0pt
% wrapper: the article is itself an amsart article; its own class and package
% list are replaced by the pinned runtime preamble plus the named packages the
% article uses, and its own macros, theorem environments and tikz libraries
% follow verbatim (source0.tex lines 3-435, without the \usepackage lines and
% without the two algorithm2e-dependent comment-styling lines, neither of which
% is used by this unit).
\usepackage{amsthm}
\usepackage{bm}
\usepackage{tikz}
\usepackage{ifthen}
\usepackage{caption}
\usepackage{subcaption}
\usepackage{wrapfig}
\usepackage{xcolor}
\usepackage{mathrsfs}
\usepackage{thmtools,thm-restate}
\providecommand{\newblock}{\hskip .11em\@plus.33em\@minus.07em}


\newcommand{\aynote}[1]{\textit{\color{red}[#1]}}

\usetikzlibrary{decorations.pathreplacing,calligraphy}
\usetikzlibrary{calc,patterns}
\newenvironment{widequote}{\begin{list}{}
      {\setlength{\rightmargin}{0mm}\setlength{\leftmargin}{5mm}}
      \item[]}{\end{list}}
\usetikzlibrary{fit}

%% as per the requirement new theorem styles can be included as shown below
%\theoremstyle{thmstyleone}%
\newtheorem{theorem}{Theorem}%  meant for continuous numbers
%%\newtheorem{theorem}{Theorem}[section]% meant for sectionwise numbers
%% optional argument [theorem] produces theorem numbering sequence instead of independent numbers for Proposition
\newtheorem{proposition}[theorem]{Proposition}% 
%%\newtheorem{proposition}{Proposition}% to get separate numbers for theorem and proposition etc.

%\theoremstyle{thmstyletwo}%
\newtheorem{example}{Example}%
\newtheorem{remark}{Remark}%

%\theoremstyle{thmstylethree}%
\newtheorem{definition}{Definition}%
\raggedbottom
%%\unnumbered% uncomment this for unnumbered level heads

\newcommand{\lsp}{\vspace{3mm}}

\newcommand{\pvct}[1]{\bm{#1}}
\newcommand{\vct}[1]{\bm{\mathsf{#1}}}
\newcommand{\mtx}[1]{\bm{\mathsf{#1}}}

\newcommand{\tikzcuboid}[5]{% width, height, depth, scale, buffer
\begin{tikzpicture}[scale=#4]
\foreach \x in {0,...,#1}
{   \draw (\x ,0  ,#3 ) -- (\x ,#2 ,#3 );
    \draw (\x ,#2 ,#3 ) -- (\x ,#2 ,0  );
}
\foreach \x in {0,...,#2}
{   \draw (#1 ,\x ,#3 ) -- (#1 ,\x ,0  );
    \draw (0  ,\x ,#3 ) -- (#1 ,\x ,#3 );
}
\foreach \x in {0,...,#3}
{   \draw (#1 ,0  ,\x ) -- (#1 ,#2 ,\x );
    \draw (0  ,#2 ,\x ) -- (#1 ,#2 ,\x );
}

\foreach \x in {0,...,#1}{
    \pgfmathtruncatemacro{\j}{mod(\x,#5)}
    \foreach \y in {0,...,#2}{
        \ifthenelse{\equal{\j}{0}}
        {\filldraw[red] (\x,\y,#3) circle (4pt);}
        {\filldraw[blue] (\x,\y,#3) circle (4pt);}
    }
}
\foreach \y in {0,...,#2}{
    \foreach \z in {0,...,#3}{
        \filldraw[red] (#1,\y,\z) circle (4pt);
    }
}
\foreach \x in {0,...,#1}{
    \pgfmathtruncatemacro{\j}{mod(\x,#5)}
    \foreach \z in {0,...,#3}{
        \ifthenelse{\equal{\j}{0}}
        {\filldraw[red] (\x,#2,\z) circle (4pt);}
        {\filldraw[blue] (\x,#2,\z) circle (4pt);}
    }
}
\pgfmathtruncatemacro{\numbuf}{#1/#5}
\pgfmathsetmacro{\halfbuf}{#5/2}
\foreach \x in {0,...,\numbuf}{
    \pgfmathtruncatemacro{\odd}{2*\x+1}
    \node [below,red] at (\x*#5,0,#3) {$I_{\odd}$};
}
\pgfmathtruncatemacro{\numbuf}{#1/#5 - 1}
\pgfmathsetmacro{\halfbuf}{#5/2}
\foreach \x in {0,...,\numbuf}{
    \pgfmathtruncatemacro{\even}{2*\x+2}
    \node [below,blue] at (\x*#5+\halfbuf,0,#3) {$I_{\even}$};
}

\draw [pen colour={black},
    decorate,
    thick,
    decoration = {calligraphic brace,
        raise=5pt,
        amplitude=5pt,
        aspect=0.5,mirror}] (#1,0) --  (#1,#2)
node[pos=0.5,right=10pt,black]{$n$};

\draw [pen colour={black},
    decorate,
    thick,
    decoration = {calligraphic brace,
        raise=5pt,
        amplitude=5pt,
        aspect=0.5}] (#1+0.5-#5,#2) --  (#1-0.5,#2)
node[pos=0.5,above=10pt,black]{$b$};
\end{tikzpicture}
}

\newcommand{\reducedtikzcuboid}[5]{% width, height, depth, scale, buffer
\begin{tikzpicture}[scale=#4]
\foreach \x in {0,...,#1}
{   \draw (\x ,0  ,#3 ) -- (\x ,#2 ,#3 );
    \draw (\x ,#2 ,#3 ) -- (\x ,#2 ,0  );
}
\foreach \x in {0,...,#2}
{   \draw (#1 ,\x ,#3 ) -- (#1 ,\x ,0  );
    \draw (0  ,\x ,#3 ) -- (#1 ,\x ,#3 );
}
\foreach \x in {0,...,#3}
{   \draw (#1 ,0  ,\x ) -- (#1 ,#2 ,\x );
    \draw (0  ,#2 ,\x ) -- (#1 ,#2 ,\x );
}

\foreach \x in {0,...,#1}{
    \pgfmathtruncatemacro{\j}{mod(\x,#5)}
    \foreach \y in {0,...,#2}{
        \ifthenelse{\equal{\j}{0}}
        {\filldraw[red] (\x,\y,#3) circle (4pt);}
        {\filldraw[gray] (\x,\y,#3) circle (4pt);}
    }
}
\foreach \y in {0,...,#2}{
    \foreach \z in {0,...,#3}{
        \filldraw[red] (#1,\y,\z) circle (4pt);
    }
}
\foreach \x in {0,...,#1}{
    \pgfmathtruncatemacro{\j}{mod(\x,#5)}
    \foreach \z in {0,...,#3}{
        \ifthenelse{\equal{\j}{0}}
        {\filldraw[red] (\x,#2,\z) circle (4pt);}
        {\filldraw[gray] (\x,#2,\z) circle (4pt);}
    }
}
\pgfmathtruncatemacro{\numbuf}{#1/#5}
\pgfmathsetmacro{\halfbuf}{#5/2}
\foreach \x in {0,...,\numbuf}{
    \pgfmathtruncatemacro{\odd}{2*\x+1}
    \node [below,red] at (\x*#5,0,#3) {$I_{\odd}$};
}

\draw [pen colour={black},
    decorate,
    thick,
    decoration = {calligraphic brace,
        raise=5pt,
        amplitude=5pt,
        aspect=0.5,mirror}] (#1,0) --  (#1,#2)
node[pos=0.5,right=10pt,black]{$n$};

\draw [pen colour={black},
    decorate,
    thick,
    decoration = {calligraphic brace,
        raise=5pt,
        amplitude=5pt,
        aspect=0.5}] (#1+0.5-#5,#2) --  (#1-0.5,#2)
node[pos=0.5,above=10pt,black]{$b$};
\end{tikzpicture}
}

\newcommand{\sparseA}[4]{% buffer, nsq, pans, scale
    \begin{tikzpicture}[scale=#4]

    \pgfmathsetmacro{\b}{#2}
    \pgfmathsetmacro{\nsq}{#1}
    \pgfmathsetmacro{\pans}{#3}
    \pgfmathsetmacro{\tot}{\nsq * (\pans + 1) + \b * \nsq * \pans}
        \draw[black](-\nsq,+\nsq) rectangle (\tot-\nsq,\tot+\nsq);
    \foreach \i in {1,...,\pans}{
        \pgfmathsetmacro{\startd}{(\i-1) * \b * \nsq + (\i-1) * \nsq}
        \pgfmathsetmacro{\endd}{\i * \b * \nsq + (\i-1) * \nsq}
        \pgfmathtruncatemacro{\l}{2*(\i)}
        \pgfmathtruncatemacro{\lpl}{2*(\i)+1}
        \pgfmathtruncatemacro{\lmn}{2*(\i)-1}
        \draw[blue] (\startd,\tot - \startd) rectangle (\endd,\tot - \endd) node[pos=.5] {$\mtx A_{\l\l}$}; % diagonal element
        \draw[blue] (\startd, \tot - \startd) rectangle (\endd, \tot - \startd + \nsq) node[pos=.5, scale=0.5] {$\mtx A_{\lmn\l}$}; % A_LC
        \draw[blue] (\startd, \tot - \startd - \nsq*\b - \nsq) rectangle (\endd, \tot - \startd + \nsq - \nsq*\b - \nsq) node[pos=.5,scale=0.5] {$\mtx A_{\lpl\l}$}; % A_RC
        \draw[blue] (\startd - \nsq, \tot - \startd) rectangle (\endd - \b * \nsq, \tot - \endd) node [pos=.5,scale=0.5] {$\mtx A_{\l\lmn}$}; % A_CL
        \draw[blue] (\startd  + \b * \nsq, \tot - \startd) rectangle (\endd + \nsq, \tot - \endd) node [pos=.5, scale=0.5] {$\mtx A_{\l\lpl}$}; % A_CL
    }
    \pgfmathsetmacro{\panspl}{\pans+1}
    \foreach \i in {1,...,\panspl}{
        \pgfmathsetmacro{\startd}{(\i-1) * \b * \nsq + (\i-2) * \nsq}
        \pgfmathsetmacro{\endd}{(\i-1) * \b * \nsq + (\i-1) * \nsq}
        \pgfmathtruncatemacro{\l}{2*(\i-1) + 1}
        \draw[red] (\startd,\tot- \startd) rectangle (\endd,\tot - \endd) node [pos = .5, scale=0.5] {$\mtx A_{\l\l}$}; % diagonal element
    }
    \pgfmathsetmacro{\tmp}{0.5*\b}
    \node [below,black] at (\tmp * \nsq * \pans+ \nsq,\nsq) {$N \times N$};
    \end{tikzpicture}
}

\newcommand{\denseT}[3]{%nsq,pans,scale
    \begin{tikzpicture}[scale=#3]
    \pgfmathsetmacro{\b}{1}
    \pgfmathsetmacro{\nsq}{#1}
    \pgfmathsetmacro{\pans}{#2}
    \pgfmathsetmacro{\tot}{\nsq * (\pans + 1) + \b * \nsq * \pans}
        \draw[black](-\nsq,+\nsq) rectangle (\tot-\nsq,\tot+\nsq);
    \foreach \i in {1,...,\pans}{
        \pgfmathsetmacro{\startd}{(\i-1) * \nsq + (\i-1) * \nsq}
        \pgfmathsetmacro{\endd}{\i * \nsq + (\i-1) * \nsq}
        \pgfmathtruncatemacro{\l}{4*(\i-1)+3}
        \pgfmathtruncatemacro{\lpl}{4*(\i)+1}
        \pgfmathtruncatemacro{\lmn}{4*(\i)-3}
        \draw[fill=red, fill opacity=0.5,text opacity=1] (\startd,\tot - \startd) rectangle (\endd,\tot - \endd) node[pos=.5,scale=0.6] {$\mtx T_{\l\l}$}; % diagonal element
        \draw[fill=red, fill opacity=0.5,text opacity=1] (\startd, \tot - \startd) rectangle (\endd, \tot - \startd + \nsq) node[pos=.5, scale=0.6] {$\mtx T_{\lmn\l}$}; % A_LC
        \draw[fill=red, fill opacity=0.5,text opacity=1] (\startd, \tot - \startd - \nsq*\b - \nsq) rectangle (\endd, \tot - \startd + \nsq - \nsq*\b - \nsq) node[pos=.5,scale=0.6] {$\mtx T_{\lpl\l}$}; % A_RC
        \draw[fill=red, fill opacity=0.5,text opacity=1] (\startd - \nsq, \tot - \startd) rectangle (\endd - \b * \nsq, \tot - \endd) node [pos=.5,scale=0.6] {$\mtx T_{\l\lmn}$}; % A_CL
        \draw[fill=red, fill opacity=0.5,text opacity=1] (\startd  + \b * \nsq, \tot - \startd) rectangle (\endd + \nsq, \tot - \endd) node [pos=.5, scale=0.6] {$\mtx T_{\l\lpl}$}; % A_CL
    }
    \pgfmathsetmacro{\panspl}{\pans+1}
    \foreach \i in {1,...,\panspl}{
        \pgfmathsetmacro{\startd}{(\i-1) * \b * \nsq + (\i-2) * \nsq}
        \pgfmathsetmacro{\endd}{(\i-1) * \b * \nsq + (\i-1) * \nsq}
        \pgfmathtruncatemacro{\l}{4*(\i-1) + 1}
        \draw[fill=red, fill opacity=0.5,text opacity=1] (\startd,\tot- \startd) rectangle (\endd,\tot - \endd) node [pos = .5, scale=0.6] {$\mtx T_{\l\l}$}; % diagonal element
    }
    \node [below,black] at (\nsq*\pans-\nsq/2,\nsq) {$\frac N b \times \frac N b$};
    \end{tikzpicture}
}

\newcommand{\squaroid}[2]{% n, scale
\begin{tikzpicture}[scale=#2]
\pgfmathsetmacro{\sep}{#1}
\pgfmathsetmacro{\seplev}{#1/2}
\pgfmathsetmacro{\n}{2*#1}

\foreach \x in {0,...,\n}
   \draw (\x ,0 ) -- (\x ,\n );

\foreach \x in {0,...,\n}
    \draw (0  ,\x ) -- (\n ,\x );


\foreach \x in {0,...,\n}
{
    \foreach \y in {0,...,\n}
    \filldraw[black!40] (\x,\y) circle (4pt);
}

\foreach \x in {0,...,\n}
{
    \filldraw[black!70] (\x,\seplev) circle (4pt);
    \filldraw[black!70] (\seplev,\x) circle (4pt);
    \filldraw[black!70] (\x,\seplev+\sep) circle (4pt);
    \filldraw[black!70] (\seplev+\sep,\x) circle (4pt);
}

\foreach \x in {0,...,\n}
{
    \filldraw[blue] (\x,\sep) circle (4pt);
    \filldraw[red] (\sep,\x) circle (4pt);
}

\node [below,red] at (\sep,0) {$I_1$};
\node [left,blue] at (0,\sep) {$I_2$};
\node [right,blue] at (\n,\sep) {$I_3$};
\node [below left,black] at (0,\n) {$I_4$};
\node [above right,black] at (\n,0) {$I_6$};
\node [above left,black] at (0,0) {$I_5$};
\node [below right,black] at (\n,\n) {$I_7$};
\end{tikzpicture}
}

\newcommand{\simpleproofpicture}[3]{ %n,b,scale
\begin{tikzpicture}[scale=#3]
\pgfmathsetmacro{\b}{#2}
\pgfmathsetmacro{\n}{#1}

\foreach \x in {0,...,\b}
   \draw (\x ,0 ) -- (\x ,\n );

\foreach \x in {0,...,\n}
    \draw (0  ,\x ) -- (\b ,\x );

\foreach \x in {0,...,\b}
{
    \foreach \y in {0,...,\n}
    \filldraw[gray] (\x,\y) circle (4pt);
}

\foreach \y in {0,...,\n}
    \filldraw[black] (0,\y) circle (4pt);

\pgfmathsetmacro{\tmpstart}{2*\b}
\pgfmathsetmacro{\tmpend}{4*\b}
\foreach \y in {\tmpstart,...,\tmpend}
    \filldraw[red] (0,\y) circle (4pt);

\node [left,black] at (0,0) {$I_1$};
\draw [pen colour={black},
    decorate,
    thick,
    decoration = {calligraphic brace,
        raise=5pt,
        amplitude=5pt,
        aspect=0.5}] (0,2*\b) --  (0,4*\b)
node[pos=0.5,left=10pt,black]{$J_B$};

\draw [pen colour={black},
    decorate,
    thick,
    decoration = {calligraphic brace,
        raise=5pt,
        amplitude=5pt,
        aspect=0.5,mirror}] (0,0) --  (\b,0)
node[pos=0.5,below=10pt,black]{$b$};
\end{tikzpicture}
}

\newcommand{\detailedproofpicture}[3]{% n,b,scale
\begin{tikzpicture}[scale=#3]
\pgfmathsetmacro{\b}{#2}
\pgfmathsetmacro{\n}{#1}

\foreach \x in {0,...,\b}
   \draw (\x ,0 ) -- (\x ,\n );

\foreach \x in {0,...,\n}
    \draw (0  ,\x ) -- (\b ,\x );

\foreach \x in {1,...,\b}
{
    \foreach \y in {0,...,\n}
    \filldraw[gray] (\x,\y) circle (4pt);
}


\foreach \y in {0,...,\n}
    \filldraw[black] (0,\y) circle (4pt);

\pgfmathsetmacro{\tmpstart}{2*\b}
\pgfmathsetmacro{\tmpend}{4*\b}
\pgfmathsetmacro{\tmpn}{\n-1}
\foreach \y in {\tmpstart,...,\tmpend}
    \filldraw[red] (0,\y) circle (4pt);


\foreach \x in {1,...,\b}
{
    \filldraw[blue] (\x,\tmpstart) circle (4pt);
    \filldraw[blue] (\x,\tmpend) circle (4pt);
}

\node [left,black] at (0,0) {$I_1$};
\node [right,black] at (\b,\tmpstart) {$J_{\gamma}$};
\node [right,black] at (\b,\tmpend) {$J_{\gamma}$};

\draw [pen colour={black},
    decorate,
    thick,
    decoration = {calligraphic brace,
        raise=5pt,
        amplitude=5pt, mirror,
        aspect=0.5}] (\b,0) --  (\b,\tmpstart-1)
node[pos=0.5,right=10pt,black]{$J_{\alpha}$};

\draw [pen colour={black},
    decorate,
    thick,
    decoration = {calligraphic brace,
        raise=5pt,
        amplitude=5pt, mirror,
        aspect=0.5}] (\b,\tmpstart+1) --  (\b,\tmpend-1)
node[pos=0.5,right=10pt,black]{$J_{\beta}$};

\draw [pen colour={black},
    decorate,
    thick,
    decoration = {calligraphic brace,
        raise=5pt,
        amplitude=5pt, mirror,
        aspect=0.5}] (\b,\tmpend+1) --  (\b,\n)
node[pos=0.5,right=10pt,black]{$J_{\alpha}$};

\draw [pen colour={black},
    decorate,
    thick,
    decoration = {calligraphic brace,
        raise=5pt,
        amplitude=5pt,
        aspect=0.5}] (0,\tmpstart) --  (0,\tmpend)
node[pos=0.5,left=10pt,black]{$J_B$};

\draw [pen colour={black},
    decorate,
    thick,
    decoration = {calligraphic brace,
        raise=5pt,
        amplitude=5pt,
        aspect=0.5,mirror}] (0.5,0) --  (\b-0.5,0)
node[pos=0.5,below=10pt,black]{$b$};
\end{tikzpicture}
}
\title[SlabLU: A Two-Level Sparse Direct Solver for Elliptic PDEs]{SlabLU: A Two-Level Sparse Direct Solver for Elliptic PDEs: the Rank Property of Thin Slabs (Proposition 1)}
\author{Anna Yesypenko and Per-Gunnar Martinsson}
\date{Source: arXiv:2211.07572v3 (CC BY 4.0)}
\begin{document}
\maketitle
\noindent\emph{Source and licence.} SlabLU: A Two-Level Sparse Direct Solver for Elliptic PDEs, arXiv:2211.07572v3, distributed under the Creative Commons Attribution 4.0 International licence, http://creativecommons.org/licenses/by/4.0/, as recorded in the arXiv OAI metadata for this exact version and shown on the abstract page https://arxiv.org/abs/2211.07572v3. Full rights notice: licenses/arxiv-2211.07572-CC-BY-4.0.txt.\par\smallskip
\noindent\textit{Editorial scope.} Counted unit: the paper\textquotesingle s Rank Property, source label \textbackslash{}label\{prop:rank\_property\} (source0.tex lines 1083--1089, printed as Proposition 1 in the article and as Proposition 1 here), delivered with its complete original proof from Appendix A at source lines 1712--1762 as the single primary proof body, including the author\textquotesingle s restatement \textbackslash{}rankprop* and the appendix figure with its caption and label. The article is itself an amsart article, so no publisher class replacement is needed: the wrapper supplies the pinned runtime preamble plus the named packages the article uses, and then the article\textquotesingle s own macro block verbatim, so that every printed number of the delivered unit (section numbers 1--3, appendix A, Proposition 1, equations (1)--(16) and (29)--(32), and the retained figures) coincides with the published version: the equations of the appendix restart from the article\textquotesingle s own value (29) through one editorial counter command, because the equations of sections 3.3 to 6 lie outside this unit. Required context retained uncounted: the whole Introduction and the whole of section 2 Discretization and node ordering with its five-point-stencil model problem, clustering and slab ordering, and the whole of sections 3.1 Schur complements and 3.2 Rank structure in the reduced blocks, that is every definition, equation, figure, remark and bibliography citation on which the counted statement and proof depend. Not part of this unit: the floating wrapfigure block of source lines 747--776, that is the illustrative five-point-stencil picture together with its caption and label. That block is left out because its own author-written tikz node helper ends every node with an extra path terminator that TeX typesets as a stray character in nullfont, which the delivery gate rejects as a missing glyph; the stencil itself is fully described by the retained text and by equation (4) at source line 725, and nothing mathematical is lost. Its single cross-reference at source line 739 is retained as a bound number. Consequently the figures of the delivered unit are numbered sequentially over the figures it actually retains (1--6), while the counts of the article\textquotesingle s own figure environment are not reproduced. Reference adaptation: the three cross-references at source lines 697, 707 and 739 point at sections 7 and 6 and at Figure 2 of the article, which lie outside this unit; their printed numbers are retained and bound to the version PDF through the reference-only mechanism, so no cross-reference is left dangling and no mathematics is touched. Citations are the article\textquotesingle s own: every \textbackslash{}cite key appearing in the retained spans is transcribed verbatim from the article\textquotesingle s own main\_v3.bbl with its published bibliography number. The article\textquotesingle s three numerical-experiment tables, its fourteen external figure assets, its algorithms and its HPS complexity sections are not part of this unit and are not redistributed. No mathematics is written, repaired or re-derived.
\section{Introduction}

\subsection{Problem setup}
We present a direct solver for boundary value problem of the form
\begin{equation}
\label{eq:bvp}
\left\{\begin{aligned}
\mathcal A u(x) =&\ f(x),\qquad&x \in \Omega,\\
u(x) =&\ g(x),\qquad&x \in \partial \Omega,
\end{aligned}\right.
\end{equation}
where $\mathcal A$ is a second order elliptic differential operator,
and $\Omega$ is a domain in two dimensions with boundary $\partial \Omega$.
The method works for a broad range of constant and variable coefficient differential
operators, but is particularly competitive for problems with highly oscillatory solutions
that are difficult to pre-condition.
For the sake of concreteness, we will focus on the case where $\mathcal A$ is a variable
coefficient Helmholtz operator
\begin{equation}
\label{eq:var_helm}
\mathcal A u(x) = -\Delta u(x) - \kappa^{2}b(x)u(x),
\end{equation}
where $\kappa$ is a reference (``typical'') wavenumber,
and where $b(x)$ is a smooth non-negative function.
Upon discretizing (\ref{eq:bvp}), one obtains a linear system
\begin{equation}
\label{eq:Au=f}
\mtx{A}\vct{u} = \vct{f},
\end{equation}
involving a coefficient matrix $\mtx{A}$ that is typically sparse. 
Our focus is on efficient algorithms for directly
building an invertible factorization of the matrix $\mtx{A}$. We specifically consider
two different discretization schemes, first a basic finite difference scheme with
second order convergence, and then a high (say $p=20$) order multidomain 
spectral 
collocation scheme \cite[Ch.~25]{2019_martinsson_book}. 
However, the techniques presented can easily be used with other (local)
discretization schemes such as finite element methods.

\subsection{Overview of proposed solver}
\label{sec:intro_overview}
The solver presented is based on a decomposition of the 
computational domain
into thin ``slabs'', as illustrated in Figure \ref{fig:domain}.
Unlike previously proposed sweeping schemes
\cite{2011_engquist_ying,gander2017restrictions,2013_sweep_demanet}
designed for preconditioning,
our objective is to directly factorize the coefficient matrix,
or at least compute a factorization that is sufficiently accurate
that it can handle problems involving strong backscattering.

To describe how the solver works, let us consider a simple model problem
where the PDE is discretized using a standard five-point finite difference
stencil on a uniform grid such as the one shown in Figure \ref{fig:domain}.
The nodes in the grid are arranged into slabs of width $b$, and are 
ordered as shown in Figure \ref{fig:domain}, resulting in a coefficient
matrix $\mtx{A}$ with the block diagonal sparsity pattern shown in 
Figure \ref{fig:sparseA}. The factorization of $\mtx{A}$ then proceeds through two stages.

\begin{figure}[!htb]
\captionsetup[subfigure]{labelfont=rm}
\centering
\begin{subfigure}[b]{0.45 \textwidth}
\centering
\tikzcuboid{16}{7}{0}{0.3}{4}
\caption{Original grid, partitioned into slabs
of width $b+2$, where $b = 3$.\\ \hspace{1em}}
\label{fig:domain}
\end{subfigure}%
\hfill
\begin{subfigure}[b]{0.45 \textwidth}
\centering
\reducedtikzcuboid{16}{7}{0}{0.3}{4}
\caption{Reduced grid, after eliminating blue nodes.
Only the red nodes are ``active''.\\}
\label{fig:reduced_grid}
\end{subfigure}

\begin{subfigure}[b]{0.45\textwidth}
    \centering
    \sparseA{18}{2.5}{4}{0.02}
    \caption{Sparsity structure of $\mtx A$ corresponding to the original grid. 
    Each block of $\mtx{A}$ is sparse.}
    \label{fig:sparseA}
\end{subfigure}%
\hfill
\begin{subfigure}[b]{0.45\textwidth}
    \centering
    \denseT{18}{2}{0.03}
    \vspace{2em}
    \caption{Sparsity structure of $\mtx T$ corresponding to the reduced grid. 
    Each block of $\mtx T$ is dense but has internal structure.}
    \label{fig:denseT}
\end{subfigure}
\caption{Illustration of the elimination order used in SlabLU.}
\end{figure}

In the first stage, the nodes that are internal to each slab 
(identified by the index vectors $I_2, I_4, \dots$ and shown as 
blue in Figure \ref{fig:domain}) are eliminated from the linear system, resulting in the reduced problem shown in
Figure \ref{fig:reduced_grid}, with the associated coefficient matrix shown
in Figure \ref{fig:denseT}. 
In this elimination step, we exploit that each subdomain is thin, which 
means that classical sparse direct solvers are particularly fast. 
To further accelerate this step, we use that the Schur complements that arise upon
the elimination of the interior nodes are rank structured. 
Specifically, they are ``HBS/HSS matrices'' \cite{2012_martinsson_FDS_survey,2013_efficient_FDS,2010_xia}
with \textit{exact} HBS/HSS rank at most $2b$. 
This allows us to accelerate this reduction step using a recently proposed 
randomized algorithm for 
compressing rank structured matrices \cite{levitt2024linear}.

The second stage is to factorize the remaining coefficient matrix $\mtx{T}$
shown in Figure \ref{fig:denseT}. This matrix is much smaller than the original
matrix $\mtx{A}$, but the sub-blocks are dense. Because we have efficiently formed
$\mtx T$ with $\mathcal H^2$-matrix structure, the reduced system can be
factorized in linear time for many elliptic PDEs (e.g. any coercive elliptic PDE, the steady state Stokes equation, and Helmholtz in the regime that the wavenumber is fixed as $N$ grows). 
With the use of $\mathcal H^2$-matrix algebra to factorize $\mtx T$, SlabLU requires linear time to 
store and factorize when the slab widths are chosen to be $\mathcal O(1)$.
In this work, we choose the slab width $b$ to grow slowly with the discretization size $N$
for performance considerations.

Our two-level framework offers a distinct advantage in terms of simplicity in both implementation and optimization. In the first stage,
we can leverage existing sparse direct solvers, which prove to be highly efficient, especially for thin 2D slabs. 
As we progress to the second stage, the fronts become larger in size, which may present a challenge when using traditional techniques.
The key benefit of SlabLU, in contrast to multi-level schemes, lies in the fact that we only need to develop specialized linear 
algebraic techniques for fronts of approximately the same size in the second stage. 
This stands in contrast to the necessity of developing such techniques for a hierarchy of front sizes in multi-level schemes.

We have found that for 2D problems, the dense 
operations are fast enough that exploiting rank structure to factorize $\mtx T$ is not
worthwhile when $N \leq 10^8$. 
Specifically, by choosing $b \sim \mathcal O(\sqrt{n})$, sparsity alone results in complexity $\mathcal O(N^{1.75})$ for the factorization stage, and $\mathcal O(N^{1.25})$ for the solve stage, when applied to 2D problems.
The simplicity of the two-level scheme allowed for parts of the factorization to 
be offloaded and accelerated on the GPU.
For meshes with 100 million points, the factorization can
be computed in 20 minutes on a desktop with an Intel i9-12900k CPU with 16 cores and an RTX 3090 GPU.
Once the factorization is available, subsequent solves take about a minute.
The numerical results feature timing results on a variety of architectures
to demonstrate that the scheme is portable to many hardware settings.

The scheme also interacts very well with high order discretization schemes such as those
described in \cite{2012_spectralcomposite} and \cite[Ch.~25]{2019_martinsson_book},
which makes it a particularly powerful tool for solving problems with highly oscillatory solutions.
The numerical results feature constant and variable coefficient 
Helmholtz problems on rectangular and curved domains. Using high order discretizations, we are able to 
discretize the PDE to 10 points per wavelength and accurately resolve
solutions on domains of size
$1000\lambda \times 1000\lambda$, where $\lambda$ is the wavelength,
to 6 digits of relative accuracy, compared to the true solution of the PDE. 


\subsection{Context and related work}

Methods to solve (\ref{eq:Au=f}) can be characterized into two groups -- direct and iterative. 
The linear systems involved are typically ill-conditioned, which necessitates the use of specialized solvers. 
For problems with non-oscillatory solutions, multigrid methods are often highly effective
\cite{mccormick_multigrid,1987_ruge_algebraic,2017_xu_algebraic}.
For oscillatory problems, multigrid works less well \cite{ernst2012difficult}. 
Specialized preconditioners have been developed, and work well for many classes of problems, in particular those involving free-space problems \cite{gander2019class,gander2007optimized,erlangga2004class,erlangga2006comparison}.
In this context, the sweeping preconditioners of Engquist and Ying are of particular relevance, as they were an inspiration for the current work \cite{2011_engquist_ying,2011_engquist_ying_PML}.
However, oscillatory problems remain highly challenging to pre-condition, in particular in situations involving strong back-scattering, cavities, or problems trapped inside a finite domain.
Sparse direct methods, which factorize the matrix $\mtx A$ exactly, offer a robust means of solving challenging PDEs. 
They are also particularly advantageous in situations involving multiple right-hand sides or low-rank updates to the matrix $\mtx A$.

The solver we describe in this work is related to multi-frontal LU solvers \cite{2016_acta_sparse_direct_survey}
which often use a hierarchical nested dissection ordering of grid nodes \cite{amestoy1996approximate,george_1973}.
For a 2D grid with $N$ nodes,
the resulting techniques have complexity $\mathcal O(N^{3/2})$ to build and $\mathcal O(N\log N)$ complexity to solve, which is known
to be work optimal among solvers that exploit only sparsity in the system 
\cite{2006_davis_directsolverbook,1989_directbook_duff}. 
The $\mathcal H$ and $\mathcal H^2$-matrix algebras can be used to reduce the complexity of operations on dense matrices that arise in 
many contexts involving the discretization of integral equations and of PDEs \cite{2008_bebendorf_book,2003_borm_introduction_H_matrix,hackbusch2015hierarchical}. 
SlabLU is inspired by prior work in sparse direct solvers for PDEs which uses $\mathcal H$-matrices \cite{amestoy2017complexity,2018_chavez_accelerated,2016_ghysels_randomized_multifrontal,2009_xia_superfast,gillman2014direct,2017_pichon_sparse}. In particular, we were inspired by \cite{2018_chavez_accelerated}
which used a domain partitioning into planes (e.g. $b=1$) for a highly effective linear solver for low order discretizations.

A key feature of SlabLU is that unlike prior work, the rank structures that we exploit are \textit{exact}, relying only on the sparsity pattern of the original coefficient matrix (cf.~Section \ref{sec:rank_reduced}).
This makes the randomized compression particularly efficient, achieving very high computational efficiency with no loss of accuracy beyond floating point errors. 
Another novelty is the usage of a recently developed black box randomized algorithm for compressing rank structured matrices \cite{levitt2022linear} (which in turn draws on insights from \cite{2011_lin_lu_ying,2016_martinsson_hudson2,2011_martinsson_randomhudson}) when eliminating the interior nodes in each slab.

Importantly, the rank-deficiencies used for thin subdomains in SlabLU are present in both the non-oscillatory and oscillatory regimes.
Similar observations are used to develop
efficient solvers for integral equation discretizations on elongated domains \cite{1996_mich_elongated,2007_martinsson_elongated}.
In the general case, the interaction rank grows algebraically with the
wavenumber \cite{engquist2018approximate}, making the efficient use of $\mathcal H$ and $\mathcal H^2$ matrix techniques challenging, though it has been observed to work well in some situations \cite{banjai2008hierarchical,betcke2017computationally,2011_xia_dehoop}. For the purposes of this work which focuses on 2D domains, 
we only use $\mathcal H^2$-matrix compression to form the reduced system $\mtx T$ efficiently,
then use highly efficient dense linear algebra routines to factorize the reduced system, an approach which we demonstrate to be effective
in the numerical results section.


%To resolve the oscillatory solutions of the Helmholtz equation, the equation must be discretized to a constant number of points per wavelength, roughly 10. PDE formulations suffer from the so-called ``pollution effect'' \cite{1997_babuska_pollution} --- as the wavenumber increases, one may suffer from phase errors unless you discretize with \textit{even more} points per wavelength. To combat these effects, we use a multi-domain high-order spectral collocation scheme, called the ``Hierarchical Poincar\'e-Steklov (HPS)'' scheme, described in \cite{babb2018accelerated,2019_martinsson_book,hao2016direct}.

%\begin{remark}
%Fast direct solvers for the Helmholtz equation that attain acceleration by exploiting 
%rank structure \pgmnote{Citations.} generally 
%lose efficiency in the high frequency regime, as the numerical 
%ranks grow with the wavenumber. 
%In contrast, the rank structures that we exploit in the first stage rely on exact algebraic
%properties, which means that they work for all wavenumbers.
%\end{remark}

\subsection{Extensions and limitations}

The solver presented is purely algebraic and can be applied to a range of different discretization schemes, including finite element and finite volume methods. 
In this manuscript, we restrict attention to regular discretizations of domains that are either rectangular themselves, 
or can be mapped smoothly to a union of elongated rectangles or slabs.
It is possible to adapt the method to more general discretizations with local refinement,
so long as it is simple to partition the computational domain into index sets corresponding to elongated slabs.
For some high order discretizations (e.g. high order finite differences), widening stencils may lead to large pre-factors when using
SlabLU, though this is also a challenge for sparse direct solvers in general \cite[Ch.~20]{2019_martinsson_book}.

While we in this manuscript restrict attention to the two dimensional case, the method is designed to handle three dimensional problems as well. 
All ideas presented carry over directly, but additional complications do arise. 
The key challenge is that in three dimensions, it is no longer feasible to use dense linear algebra when factorizing the block tridiagonal reduced coefficient matrix $\mtx{T}$.
However, the 3D version of SlabLU is also very easy to parallelize, and the idea of using randomized compression combined with efficient sparse direct solvers to eliminate the nodes interior to each slab still applies, cf.~Section \href{https://arxiv.org/abs/2211.07572v3}{7}.

\section{Discretization and node ordering}

We introduce two different discretization techniques for
(\ref{eq:bvp}). The first is simply the standard second order accurate
five point finite difference stencil. Since this discretization is very 
well known, it allows us to describe how the solver works without the need
to introduce cumbersome background material. To demonstrate that the solver
works for a broader class of discretization schemes, the numerical experiments
reported in Section \href{https://arxiv.org/abs/2211.07572v3}{6} also include results that rely on the high 
order accurate \textit{Hierarchical Poincar\'e-Steklov (HPS)} scheme, which
we briefly describe in Section \ref{sec:other_disc}.

\subsection{A model problem based on the five point stencil}
\label{sec:modeldisc}

For purposes of describing the factorization scheme, let us introduce a very
simple discretization of the boundary value problem (\ref{eq:bvp}).
We work 
with a rectangular domain $\Omega=[0,L_1] \times [0,L_2]$ and the second order
linear elliptic operator $\mathcal A$ defined by (\ref{eq:var_helm}). We assume 
that $L_1 \ge L_2$, and that $L_{1} = hn_{1}$ and $L_{2} = hn_{2}$ for some grid 
spacing $h$ and some positive integers $n_{1}$ and $n_{2}$.
We then
discretize $\mathcal{A}$ with a standard second-order finite difference scheme,
to obtain the linear system
\begin{equation}
\label{eq:discretiation}
\frac 1 {h^2} 
\bigl(
\vct{u}(n_{\rm w}) + 
\vct{u}(n_{\rm e}) +
\vct{u}(n_{\rm n}) +
\vct{u}(n_{\rm s}) -
4 \vct{u}(n)
\bigr)
- \kappa^2 \vct{b}(n)\vct{u}(n)
=
\vct{f}(n).
\end{equation}
The vector $\vct{f}$ holds values of the body load at the discretization nodes, and the vector $u$ holds approximations to the solution $u$. 
See Figure \href{https://arxiv.org/abs/2211.07572v3}{2 of the published article} for a visualization of the 5 point stencil. 
We write the system (\ref{eq:discretiation}) compactly as
$\mtx{A}\vct{u} = \vct{f}$.


\subsection{Clustering of the nodes}
\label{sec:clustering}


We next subdivide the computational domain into thin ``slabs'', as shown in 
Figure \ref{fig:domain}. We let $b$ denote the number of grid points in each
slab ($b=4$ in Figure \ref{fig:domain}), and then introduce index vectors $I_{1},\, I_{2},\, I_{3}, \dots$ that keep
track of which slabs each grid point belongs to. The odd numbered index
vectors $I_{1},\, I_{3},\,I_{5},\dots$ indicate nodes on the interfaces between
slabs (red in Figure \ref{fig:domain}), while the even numbered ones indicate
nodes that are interior to each slab (blue in Figure \ref{fig:domain}). 
With this ordering of the grid points, the stiffness matrix associated with the
discretization (\ref{eq:discretiation}) has the sparsity pattern shown in 
Figure \ref{fig:sparseA}.


\subsection{High order discretizations}
\label{sec:other_disc}

\section{Stage One: Elimination of nodes interior to each slab}
\label{sec:stageone}

This section describes the process that we use to eliminate 
the nodes interior to each slab that we sketched out in Section \ref{sec:intro_overview}.
The objective is to reduce the sparse stiffness matrix $\mtx{A}$ 
(illustrated in Figure \ref{fig:sparseA}) into the smaller block
tridiagonal matrix $\mtx{T}$ (illustrated in Figure \ref{fig:denseT}).
The techniques described form the core algorithmic innovation
of the manuscript.

\subsection{Schur complements}
With the ordering introduced in Section \ref{sec:clustering}, the
coefficient matrix $\mtx{A}$ has the block form
\begin{equation}
\label{eq:dell1}
\begin{bmatrix}
\mtx A_{11} & \mtx A_{12} & \mtx 0 & \mtx 0 & \mtx 0 & \dots \\
\mtx A_{21} & \mtx A_{22} & \mtx A_{23} & \mtx 0 &  \mtx 0 & \dots \\
\mtx 0 & \mtx A_{32} & \mtx A_{33} & \mtx A_{34} &   \mtx 0 & \dots\\
\mtx{0} & \mtx 0 & \mtx A_{43} & \mtx A_{44} & \mtx A_{45} & \dots\\
\vdots & \vdots & \vdots & \vdots & \vdots & \vdots
\end{bmatrix}
\begin{bmatrix} \mtx u_1\\\mtx u_2\\ \mtx u_3 \\ \mtx u_4 \\ \vdots \end{bmatrix} =
\begin{bmatrix} \mtx f_1\\ \mtx f_2\\ \mtx f_3\\ \mtx f_4 \\ \vdots        \end{bmatrix}.
\end{equation}
We eliminate the vectors $\vct{u}_{2},\,\vct{u}_{4},\,\vct{u}_{6},\,\dots$
that represent unknown variables in the interior of each slab through a 
step of block Gaussian elimination. To be precise, we insert the relation
\begin{equation}
\vct{u}_{i} = 
\mtx{A}_{ii}^{-1}\bigl(\vct{f}_{i} - \mtx{A}_{i,i-1}\vct{u}_{i-1} - \mtx{A}_{i,i+1}\vct{u}_{i+1}\bigr),
\qquad i = 2,\,4,\,6,\,\dots
\label{eq:usol_even}
\end{equation}
into the odd-numbered rows in (\ref{eq:dell1}) to obtain the reduced system
\begin{equation}
\label{eq:dell3}
\begin{bmatrix}
\mtx T_{11} & \mtx T_{13} & \mtx 0 & \mtx 0 & \mtx 0 & \dots \\
\mtx T_{31} & \mtx T_{33} & \mtx T_{35} & \mtx 0 &  \mtx 0 & \dots \\
\mtx 0 & \mtx T_{53} & \mtx T_{55} & \mtx T_{57} &   \mtx 0 & \dots\\
\mtx{0} & \mtx 0 & \mtx T_{57} & \mtx T_{77} & \mtx T_{79} & \dots\\
\vdots & \vdots & \vdots & \vdots & \vdots & \vdots
\end{bmatrix}
\begin{bmatrix} \mtx u_1\\\mtx u_3\\ \mtx u_5 \\ \mtx u_7 \\ \vdots \end{bmatrix} =
\begin{bmatrix} \tilde{\vct{f}}_{1} \\ \tilde{\vct{f}}_{3} \\\tilde{\vct{f}}_{5} \\ \tilde{\mtx f}_7\\ \vdots \end{bmatrix},
\end{equation}
where the sub-blocks of $\mtx T$ are defined as
\begin{align}
\label{eq:T11}
\mtx{T}_{11} &= \mtx A_{11} - \mtx{A}_{ 12}\ \mtx{A}_{ 22}^{-1}\ \mtx{A}_{ 21}, \\
\label{eq:T13}
\mtx{T}_{13} &= \mtx A_{13} - \mtx{A}_{ 12}\ \mtx{A}_{ 22}^{-1}\ \mtx{A}_{ 23}, \\
\label{eq:T31}
\mtx{T}_{31} &= \mtx A_{31} - \mtx{A}_{ 32}\ \mtx{A}_{ 22}^{-1}\ \mtx{A}_{ 23}, \\
\label{eq:T33}
\mtx{T}_{33} &= \mtx A_{33} - \mtx{A}_{ 32}\ \mtx{A}_{ 22}^{-1}\ \mtx{A}_{ 23}
                              - \mtx{A}_{ 34}\ \mtx{A}_{ 44}^{-1}\ \mtx{A}_{ 43}, \\
\label{eq:T35}
\mtx{T}_{35} &= \mtx A_{35} - \mtx{A}_{ 34}\ \mtx{A}_{ 44}^{-1}\ \mtx{A}_{ 23},
\end{align}
and so on. The reduced right-hand sides $\mtx{\tilde f}$ are defined as
\begin{align}
\label{eq:f1}
\mtx{\tilde f}_{ 1} &= \mtx f_{ 1} - \mtx{A}_{ 12}\ \mtx{A}_{22}^{-1}\ \mtx f_2, \\
\label{eq:f3}
\mtx{\tilde f}_{ 3} &= \mtx f_{ 3} - \mtx{A}_{ 32}\ \mtx{A}_{22}^{-1}\ \mtx f_2 
                                   - \mtx{A}_{ 34}\ \mtx{A}_{44}^{-1}\ \mtx f_4, \\
\label{eq:f5}
\mtx{\tilde f}_{ 5} &= \mtx f_{ 5} - \mtx{A}_{ 54}\ \mtx{A}_{44}^{-1}\ \mtx f_4
                                   - \mtx{A}_{ 56}\ \mtx{A}_{66}^{-1}\ \mtx f_6,
\end{align}
etc.

\subsection{Rank structure in the reduced blocks}
\label{sec:rank_reduced}

We next discuss algebraic properties of the blocks of the reduced coefficient matrix that allow for $\mtx T$ to be formed efficiently. The sub-blocks of $\mtx T$ are Schur-complements of sparse matrices. For instance, block $\mtx T_{\rm 11}$
has the formula
\begin{equation}\underset{n_2 \times n_2}{\mtx T_{11}} = \underset{n_2 \times n_2}{\mtx A_{11}} - \underset{n_2 \times n_2 b} {\mtx A_{12}}\ 
\underset{n_2 b \times n_2 b} {\mtx A_{22}^{-1}}\ \underset{n_2 b \times n_2} {\mtx A_{21}},
\label{eq:T11_detailed}
\end{equation}
where $\mtx A_{11}, \mtx A_{12}, \mtx A_{21}$ are sparse with $\mathcal O(n_2)$ non-zero entries and $\mtx A_{22}$ is a sparse banded matrix. 
Factorizing $\mtx A_{22}$ can be done efficiently using a sparse direct solver, 
but forming $\mtx A_{22}^{-1} \mtx A_{21}$ naively may be slow and memory-intensive, especially considering that 
$\mtx A_{21}$ is sparse and would need to be converted to dense to interface with solve routines. 

We use an alternate approach for efficiently forming $\mtx T_{11}$ that achieves high arithmetic intensity
while maintaining a very low memory footprint. First, we prove algebraic properties about $\mtx T_{11}$, which is a
dense but structured matrix that only needs $\mathcal O(n_2 b)$ storage 
in exact precision. 
The matrix $\mtx T_{11}$ is compressible in a format called Hierarchically Block-Separable (HBS) or Hierarchically
Semi-Separable (HSS) with exact HBS rank at most $2b$. 
HBS matrices are a type of hierarchical matrix ($\mathcal H^2$-matrix), which allow dense matrices to be stored efficiently
by exploiting low-rank structure in sub-blocks at different levels of granularity \cite{2008_bebendorf_book,hackbusch2015hierarchical,2019_martinsson_book}.
The sub-blocks of $\mtx T$ are also compressible in $\mathcal H$-matrix formats, e.g. Hierarchical Off-Diagonal Low Rank (HODLR). 
The rank property of $\mtx T_{11}$ is formally stated in Proposition \ref{prop:rank_property}.

\begin{restatable}[Rank Property]{proposition}{rankprop}
\label{prop:rank_property}
Let $J_B$ be a contiguous set of points on the slab interface $I_1$, and 
let $J_F$ be the rest of the points $J_{F} = I_1 \setminus J_{B}$.
The sub-matrices
$(\mtx T_{11})_{BF},\ (\mtx T_{11})_{FB}$ of the matrix $\mtx T$ have exact rank at most $2b$. 
\end{restatable}

\appendix

\section{Rank Property of Thin Slabs}
\label{sec:appendix}

\setcounter{equation}{28}
In this appendix, we prove Proposition \ref{prop:rank_property}, which makes
a claim on the rank structure of $\mtx T_{11}$, defined in (\ref{eq:T11_detailed}).
\rankprop*

\begin{figure}[!htb]
\centering
\begin{minipage}{0.4\textwidth}
\centering
\detailedproofpicture{28}{5}{0.20}
\end{minipage}%
\begin{minipage}{0.4\textwidth}
\captionof{figure}{To assist in the proof of Proposition \ref{prop:rank_property}, we define a partitioning of the slab interface $I_1 = J_B \cup J_F$, where $J_F = I_1 \setminus J_B$. We also partition the slab interior nodes into $I_2 := J_{\alpha} \cup J_{\beta} \cup J_{\gamma}$.}
\label{fig:partI2}
\end{minipage}
\end{figure}

Recall that $\mtx T_{11} = \mtx A_{11} - \mtx A_{12} \mtx A_{22}^{-1} \mtx A_{21}$. The proof relies on the sparsity structure of the matrices in the Schur complement. As stated in the proposition, the slab interface $I_1$ is partitioned into
indices $J_B$ and $J_F$. The proof relies on partitioning $I_2$ as well, into the indices $J_{\alpha},J_{\beta}, J_{\gamma}$ shown in Figure \ref{fig:partI2}, where $|J_{\gamma}| = 2b$.

The matrix $\mtx A_{22}$ is sparse and can be factorized as
\begin{equation}
\mtx A_{22} = \mtx L_{22} \mtx U_{22} := 
\begin{bmatrix} \mtx L_{\alpha \alpha}\\
& \mtx L_{\beta \beta}\\
\mtx L_{\gamma \alpha} & \mtx L_{\gamma \beta} & \mtx L_{\gamma \gamma}\end{bmatrix} 
\begin{bmatrix} \mtx U_{\alpha \alpha} & & \mtx U_{\alpha \gamma}\\
& \mtx U_{\beta \beta} & \mtx U_{\beta \gamma}\\
 &  & \mtx U_{\gamma \gamma}\end{bmatrix} 
 \label{eq:sparse_A22}
\end{equation}
The formula for ${\left( \mtx T_{11} \right)}_{FB}$ can be re-written as
\begin{equation}
{\left( \mtx T_{11} \right)}_{FB} = \mtx A_{FB}
 - {\left( \mtx A_{F2} \mtx U_{22}^{-1}\right)}
 {\left( \mtx L_{22}^{-1} \mtx A_{2B}\right)}
:= \mtx A_{FB} - {\mtx X_{F2}}\ {\mtx Y_{2B}}
\end{equation}
The factors $\mtx X_{F2}$ and $\mtx Y_{2B}$ have sparse
structure, due the sparsity in the factorization (\ref{eq:sparse_A22})
and the sparsity of $\mtx A_{F2}$ and $\mtx A_{2B}$.
\begin{equation}
\mtx X_{F2} = \begin{bmatrix} \mtx A_{F\alpha}& \mtx 0 & \mtx A_{F\gamma}\end{bmatrix} \mtx U_{22}^{-1},\qquad \mtx Y_{2B} 
= \mtx L_{22}^{-1} \begin{bmatrix} \mtx 0\\ \mtx A_{\beta B}\\ \mtx A_{\gamma B}\end{bmatrix}
\end{equation}
The factors $\mtx X_{F2}$ and $\mtx Y_{2B}$ have the same sparsity pattern as $\mtx A_{F2}$ and $\mtx A_{2B}$, respectively.
As a result,
\begin{equation}
{\left( \mtx T_{11} \right)}_{FB} = \mtx A_{FB} - \begin{bmatrix} \mtx X_{F\alpha}& \mtx 0 & \mtx X_{F\gamma}\end{bmatrix} \begin{bmatrix} \mtx 0\\ \mtx Y_{\beta B}\\ \mtx Y_{\gamma B}\end{bmatrix} = \underset{\text{sparse, $\mathcal O(1)$ entries}}{\mtx A_{F,B}} - 
\underset{\text{exact rank}\ 2b}{\mtx X_{F\gamma} \mtx Y_{\gamma B}}.
\end{equation}
Similar reasoning can be used to show the result for $\left( \mtx T_{11} \right)_{BF}$.

\begin{thebibliography}{99}
\bibitem[3]{amestoy2017complexity} Patrick Amestoy, Alfredo Buttari, Jean-Yves l'Excellent, and Theo Mary.
\newblock On the complexity of the block low-rank multifrontal factorization.
\newblock {\em SIAM Journal on Scientific Computing}, 39(4):A1710--A1740, 2017.
\bibitem[4]{amestoy1996approximate} Patrick~R Amestoy, Timothy~A Davis, and Iain~S Duff.
\newblock An approximate minimum degree ordering algorithm.
\newblock {\em SIAM Journal on Matrix Analysis and Applications},
  17(4):886--905, 1996.
\bibitem[8]{banjai2008hierarchical} Lehel Banjai and Wolfgang Hackbusch.
\newblock Hierarchical matrix techniques for low-and high-frequency helmholtz
  problems.
\newblock {\em IMA journal of numerical analysis}, 28(1):46--79, 2008.
\bibitem[10]{2008_bebendorf_book} Mario Bebendorf.
\newblock {\em Hierarchical matrices}, volume~63 of {\em Lecture Notes in
  Computational Science and Engineering}.
\newblock Springer-Verlag, Berlin, 2008.
\newblock A means to efficiently solve elliptic boundary value problems.
\bibitem[12]{betcke2017computationally} Timo Betcke, Elwin van~’t Wout, and Pierre G{\'e}lat.
\newblock Computationally efficient boundary element methods for high-frequency
  helmholtz problems in unbounded domains.
\newblock {\em Modern Solvers for Helmholtz Problems}, pages 215--243, 2017.
\bibitem[14]{2003_borm_introduction_H_matrix} Steffen B{\"o}rm, Lars Grasedyck, and Wolfgang Hackbusch.
\newblock Introduction to hierarchical matrices with applications.
\newblock {\em Engineering analysis with boundary elements}, 27(5):405--422,
  2003.
\bibitem[15]{mccormick_multigrid} William~L. Briggs, Van~Emden Henson, and Steve~F. McCormick.
\newblock {\em A multigrid tutorial}.
\newblock Society for Industrial and Applied Mathematics (SIAM), Philadelphia,
  PA, second edition, 2000.
\bibitem[16]{2018_chavez_accelerated} Gustavo Ch{\'a}vez, George Turkiyyah, Stefano Zampini, Hatem Ltaief, and David
  Keyes.
\newblock Accelerated cyclic reduction: A distributed-memory fast solver for
  structured linear systems.
\newblock {\em Parallel Computing}, 74:65--83, 2018.
\bibitem[18]{2006_davis_directsolverbook} Timothy~A Davis.
\newblock {\em Direct methods for sparse linear systems}, volume~2.
\newblock {SIAM}, Philadelphia PA, 2006.
\bibitem[19]{2016_acta_sparse_direct_survey} Timothy~A. Davis, Sivasankaran Rajamanickam, and Wissam~M. Sid-Lakhdar.
\newblock A survey of direct methods for sparse linear systems.
\newblock {\em Acta Numerica}, 25:383 -- 566, 2016.
\bibitem[21]{1989_directbook_duff} I.S. Duff, A.M. Erisman, and J.K. Reid.
\newblock {\em Direct Methods for Sparse Matrices}.
\newblock Oxford, Oxford United Kingdom, 1989.
\bibitem[22]{2011_engquist_ying} B.~Engquist and L.~Ying.
\newblock Sweeping preconditioner for the {Helmholtz} equation: Hierarchical
  matrix representation.
\newblock {\em Communications on Pure and Applied Mathematics}, 64(5):697--735,
  2011.
\bibitem[23]{2011_engquist_ying_PML} Bj\"{o}rn Engquist and Lexing Ying.
\newblock Sweeping preconditioner for the {Helmholtz} equation: Moving
  perfectly matched layers.
\newblock {\em Multiscale Modeling \& Simulation}, 9(2):686--710, 2011.
\bibitem[24]{engquist2018approximate} Bj{\"o}rn Engquist and Hongkai Zhao.
\newblock Approximate separability of the green's function of the helmholtz
  equation in the high frequency limit.
\newblock {\em Communications on Pure and Applied Mathematics},
  71(11):2220--2274, 2018.
\bibitem[25]{erlangga2006comparison} Yogi~A Erlangga, Cornelis Vuik, and Cornelis~W Oosterlee.
\newblock Comparison of multigrid and incomplete lu shifted-laplace
  preconditioners for the inhomogeneous helmholtz equation.
\newblock {\em Applied numerical mathematics}, 56(5):648--666, 2006.
\bibitem[26]{erlangga2004class} Yogi~A Erlangga, Cornelis Vuik, and Cornelis~Willebrordus Oosterlee.
\newblock On a class of preconditioners for solving the helmholtz equation.
\newblock {\em Applied Numerical Mathematics}, 50(3-4):409--425, 2004.
\bibitem[27]{ernst2012difficult} Oliver~G Ernst and Martin~J Gander.
\newblock Why it is difficult to solve {Helmholtz} problems with classical
  iterative methods.
\newblock {\em Numerical analysis of multiscale problems}, pages 325--363,
  2012.
\bibitem[31]{gander2007optimized} Martin~J Gander, Laurence Halpern, and Fr{\'e}d{\'e}ric Magoules.
\newblock An optimized schwarz method with two-sided robin transmission
  conditions for the helmholtz equation.
\newblock {\em International journal for numerical methods in fluids},
  55(2):163--175, 2007.
\bibitem[32]{gander2017restrictions} Martin~J Gander and Hui Zhang.
\newblock Restrictions on the use of sweeping type preconditioners for
  {Helmholtz} problems.
\newblock In {\em International Conference on Domain Decomposition Methods},
  pages 321--332. Springer, 2017.
\bibitem[33]{gander2019class} Martin~J Gander and Hui Zhang.
\newblock A class of iterative solvers for the {Helmholtz} equation:
  Factorizations, sweeping preconditioners, source transfer, single layer
  potentials, polarized traces, and optimized schwarz methods.
\newblock {\em Siam Review}, 61(1):3--76, 2019.
\bibitem[35]{george_1973} A.~George.
\newblock Nested dissection of a regular finite element mesh.
\newblock {\em SIAM J. on Numerical Analysis}, 10:345--363, 1973.
\bibitem[36]{2016_ghysels_randomized_multifrontal} P.~Ghysels, X.~Li, F.~Rouet, S.~Williams, and A.~Napov.
\newblock {An Efficient Multicore Implementation of a Novel HSS-Structured
  Multifrontal Solver Using Randomized Sampling}.
\newblock {\em SIAM Journal on Scientific Computing}, 38(5):S358--S384, 2016.
\bibitem[40]{gillman2014direct} Adrianna Gillman and Per-Gunnar Martinsson.
\newblock A direct solver with {O(N)} complexity for variable coefficient
  elliptic {PDE}s discretized via a high-order composite spectral collocation
  method.
\newblock {\em SIAM Journal on Scientific Computing}, 36(4):A2023--A2046, 2014.
\bibitem[41]{2012_martinsson_FDS_survey} Adrianna Gillman, Patrick Young, and Per-Gunnar Martinsson.
\newblock A direct solver $o(n)$ complexity for integral equations on
  one-dimensional domains.
\newblock {\em Frontiers of Mathematics in China}, 7:217--247, 2012.
\newblock 10.1007/s11464-012-0188-3.
\bibitem[42]{hackbusch2015hierarchical} Wolfgang Hackbusch.
\newblock {\em Hierarchical matrices: algorithms and analysis}, volume~49.
\newblock Springer, New York NY, 2015.
\bibitem[45]{levitt2022linear} James Levitt and Per-Gunnar Martinsson.
\newblock {Linear-Complexity Black-Box Randomized Compression of Hierarchically
  Block Separable Matrices}.
\newblock {\em arXiv preprint arXiv:2205.02990}, 2022.
\bibitem[46]{levitt2024linear} James Levitt and Per-Gunnar Martinsson.
\newblock Linear-complexity black-box randomized compression of rank-structured
  matrices.
\newblock {\em SIAM Journal on Scientific Computing}, 46(3):A1747--A1763, 2024.
\bibitem[49]{2011_lin_lu_ying} L.~Lin, J.~Lu, and L.~Ying.
\newblock Fast construction of hierarchical matrix representation from
  matrix-vector multiplication.
\newblock {\em Journal of Computational Physics}, 230(10):4071 -- 4087, 2011.
\bibitem[51]{2016_martinsson_hudson2} Per-Gunnar Martinsson.
\newblock Compressing rank-structured matrices via randomized sampling.
\newblock {\em SIAM Journal on Scientific Computing}, 38(4):A1959--A1986, 2016.
\bibitem[52]{2019_martinsson_book} Per-Gunnar Martinsson.
\newblock {\em Fast direct solvers for elliptic PDEs}.
\newblock {SIAM}, Philadelphia PA, 2019.
\bibitem[54]{2011_martinsson_randomhudson} P.G. Martinsson.
\newblock A fast randomized algorithm for computing a hierarchically
  semiseparable representation of a matrix.
\newblock {\em SIAM Journal on Matrix Analysis and Applications},
  32(4):1251--1274, 2011.
\bibitem[55]{2012_spectralcomposite} P.G. Martinsson.
\newblock A direct solver for variable coefficient elliptic pdes discretized
  via a composite spectral collocation method.
\newblock {\em Journal of Computational Physics}, 242(0):460 -- 479, 2013.
\bibitem[56]{2007_martinsson_elongated} P.G. Martinsson and V.~Rokhlin.
\newblock A fast direct solver for scattering problems involving elongated
  structures.
\newblock {\em Journal of Computational Physics}, 221:288--302, 2007.
\bibitem[57]{1996_mich_elongated} E.~Michielssen, A.~Boag, and W.~C. Chew.
\newblock Scattering from elongated objects: direct solution in
  ${O}({N}\log^{2}{N})$ operations.
\newblock {\em IEE Proc. Microw. Antennas Propag.}, 143(4):277 -- 283, 1996.
\bibitem[59]{2017_pichon_sparse} Gr{\'e}goire Pichon, Eric Darve, Mathieu Faverge, Pierre Ramet, and Jean Roman.
\newblock Sparse supernodal solver using block low-rank compression.
\newblock In {\em 2017 IEEE International Parallel and Distributed Processing
  Symposium Workshops (IPDPSW)}, pages 1138--1147. IEEE, 2017.
\bibitem[60]{1987_ruge_algebraic} John~W Ruge and Klaus St{\"u}ben.
\newblock Algebraic multigrid.
\newblock In {\em Multigrid methods}, pages 73--130. SIAM, Philadelphia PA,
  1987.
\bibitem[61]{2013_sweep_demanet} Alexandre Vion, R~B{\'e}langer-Rioux, L~Demanet, and Christophe Geuzaine.
\newblock A {DDM} double sweep preconditioner for the {Helmholtz} equation with
  matrix probing of the {DtN} map.
\newblock {\em Mathematical and Numerical Aspects of Wave Propagation WAVES},
  2013, 2013.
\bibitem[63]{2011_xia_dehoop} Shen Wang, Maarten~V. de~Hoop, and Jianlin Xia.
\newblock On 3d modeling of seismic wave propagation via a structured parallel
  multifrontal direct {Helmholtz} solver.
\newblock {\em Geophysical Prospecting}, 59(5):857--873, 2011.
\bibitem[64]{2013_efficient_FDS} Shen Wang, Xiaoye~S Li, Jianlin Xia, Yingchong Situ, and Maarten~V De~Hoop.
\newblock Efficient scalable algorithms for solving dense linear systems with
  hierarchically semiseparable structures.
\newblock {\em SIAM Journal on Scientific Computing}, 35(6):C519--C544, 2013.
\bibitem[65]{2010_xia} J.~Xia, S.~Chandrasekaran, M.~Gu, and X.S. Li.
\newblock Fast algorithms for hierarchically semiseparable matrices.
\newblock {\em Numerical Linear Algebra with Applications}, 17(6):953--976,
  2010.
\bibitem[66]{2009_xia_superfast} Jianlin Xia, Shivkumar Chandrasekaran, Ming Gu, and Xiaoye~S. Li.
\newblock Superfast multifrontal method for large structured linear systems of
  equations.
\newblock {\em SIAM J. Matrix Anal. Appl.}, 31(3):1382--1411, 2010.
\bibitem[67]{2017_xu_algebraic} Jinchao Xu and Ludmil Zikatanov.
\newblock Algebraic multigrid methods.
\newblock {\em Acta Numerica}, 26:591--721, 2017.
\end{thebibliography}
\end{document}
